\documentclass[12pt,a4paper]{article}

% -------------------------------------------------
% Sprache / Kodierung
% -------------------------------------------------
\usepackage[T1]{fontenc}
\usepackage[utf8]{inputenc}
\usepackage[ngerman]{babel}
\usepackage{lmodern}
\usepackage{microtype}

% -------------------------------------------------
% Mathematik
% -------------------------------------------------
\usepackage{amsmath,amssymb,amsthm,mathtools}

% -------------------------------------------------
% Layout / Grafik / Farben
% -------------------------------------------------
\usepackage[a4paper,margin=2.3cm]{geometry}
\usepackage{booktabs}
\usepackage{xcolor}
\usepackage[hidelinks,unicode]{hyperref}
\pdfstringdefDisableCommands{%
  \def\ü{ue}\def\ö{oe}\def\ä{ae}\def\Ü{Ue}\def\Ö{Oe}\def\Ä{Ae}\def\ss{ss}%
}
\usepackage[most]{tcolorbox}
\usepackage{enumitem}
\usepackage{array}
\usepackage{listings}

% -------------------------------------------------
% Farben
% -------------------------------------------------
\definecolor{lückenrot}{RGB}{180,40,40}
\definecolor{bewiesengrün}{RGB}{30,100,50}
\definecolor{warnorange}{RGB}{200,100,10}
\definecolor{hintergrundblau}{RGB}{235,242,250}
\definecolor{hintergrundgelb}{RGB}{255,252,225}
\definecolor{adischviolett}{RGB}{90,20,140}
\definecolor{fehlerrot}{RGB}{200,0,0}
\definecolor{konjviolett}{RGB}{80,0,130}
\definecolor{leanblau}{RGB}{0,80,160}
\definecolor{mathlibgrün}{RGB}{20,120,50}
\definecolor{roadmaporange}{RGB}{160,80,0}
\definecolor{codeblau}{RGB}{40,60,130}

% -------------------------------------------------
% tcolorbox-Stile
% -------------------------------------------------
\tcbset{
  lücke/.style={colback=red!8, colframe=lückenrot, fonttitle=\bfseries,
    title={Fehlt in Mathlib (rot)}},
  vorhanden/.style={colback=green!7, colframe=mathlibgrün,
    fonttitle=\bfseries, title={Formalisiert in Mathlib (grün)}},
  teilweise/.style={colback=orange!8, colframe=warnorange,
    fonttitle=\bfseries, title={Teilweise formalisiert (orange)}},
  roadmap/.style={colback=yellow!8, colframe=roadmaporange,
    fonttitle=\bfseries, title={Roadmap-Lemma}},
  lean/.style={colback=blue!5, colframe=leanblau, fonttitle=\bfseries},
  kernidee/.style={colback=hintergrundblau, colframe=blue!60!black,
    fonttitle=\bfseries, title={Kernidee}},
  warnung/.style={colback=orange!10, colframe=warnorange,
    fonttitle=\bfseries, title={Hinweis}},
}

% -------------------------------------------------
% Listings für Lean-Code
% -------------------------------------------------
\lstdefinelanguage{Lean4}{
  keywords={theorem,lemma,def,noncomputable,import,open,where,by,
    exact,simp,rw,intro,apply,have,show,fun,let,sorry,structure,
    namespace,end,variable,instance,class,inductive,if,then,else,
    forall,exists,match,with,return,type,Prop,Type,Sort},
  keywordstyle=\color{leanblau}\bfseries,
  comment=[l]{--},
  commentstyle=\color{gray}\itshape,
  string=[b]{"},
  stringstyle=\color{bewiesengrün},
  basicstyle=\ttfamily\small,
  breaklines=true,
  frame=single,
  rulecolor=\color{leanblau!40},
  backgroundcolor=\color{blue!3},
}

% -------------------------------------------------
% Theorem-Umgebungen
% -------------------------------------------------
\newtheorem{definition}{Definition}[section]
\newtheorem{proposition}[definition]{Proposition}
\newtheorem{theorem}[definition]{Theorem}
\newtheorem{lemma}[definition]{Lemma}
\newtheorem{korollar}[definition]{Korollar}
\newtheorem{bemerkung}[definition]{Bemerkung}

% -------------------------------------------------
% Makros
% -------------------------------------------------
\newcommand{\vii}{\nu_2}
\newcommand{\N}{\mathbb{N}}
\newcommand{\Z}{\mathbb{Z}}
\newcommand{\R}{\mathbb{R}}
\newcommand{\Q}{\mathbb{Q}}
\newcommand{\Mathlib}[1]{\texttt{Mathlib.\allowbreak #1}}

% -------------------------------------------------
% Dokument
% -------------------------------------------------
\title{%
  Lean\,4\,/\,Mathlib-Recherche für die Collatz-Beweisstrategie\\[0.3em]
  \large Systematische Analyse formalisierter Theoreme\\[0.3em]
  \large Verfügbare Bausteine, offene Lücken und Formalisierungs-Roadmap%
}
\author{Thomas Hoffbauer}
\date{16.\ Juni 2026 \quad\textit{(Lean 4.30.0, Mathlib4)}}

\begin{document}
\maketitle

\begin{abstract}
Diese Datei dokumentiert eine systematische Mathlib4-Recherche für die Bausteine der
Collatz-Beweisstrategie aus \texttt{collatz\_lte\_dichte.tex},
\texttt{collatz\_vier\_teilnormen.tex} und den zugehörigen Begleitdokumenten.

\textbf{Kernbefund:} Stirling-Formel (inkl.\ Robbins-Schranke), Bernoulli-Zahlen und
von-Staudt-Clausen, Lagrange-Vierquadratesatz sowie die gesamte Infrastruktur für
geometrische Reihen ($x=\tfrac12$) und 2-adische Valuationen sind in Mathlib vollständig
formalisiert. Die Ergodizitätsdefinition und der von-Neumann-Mittelwertsatz sind vorhanden.

\textbf{Kernlücken:} Das Dichte-Lemma für C-Ketten-Startpunkte
($\mathrm{density}\,\{n \mid \vii(n+1)\geq k\} = 2^{-k}$), die Kingman-Subadditivität,
Collatz-spezifische Strukturlemmata und der Birkhoff-Punktgrenzwertsatz für
Wahrscheinlichkeitsmaße fehlen noch.

\textbf{Umgebung:} Lean~4.30.0 (elan) installiert; Mathlib als Bibliothek über
\texttt{lake} einbindbar. Kein Lean-Mathlib-Projekt im aktuellen Arbeitsverzeichnis.
\end{abstract}

\tableofcontents
\bigskip

\begin{tcolorbox}[warnung]
\textbf{Grundlage dieser Recherche:} Mathlib4-Onlinedokumentation
(\url{https://leanprover-community.github.io/mathlib4_docs/}),
Stand: 16.\ Juni 2026 (Lean 4.30.0).
Die Lean-Dateien im Arbeitsverzeichnis (\texttt{Bertrand.lean},
\texttt{LeanBeweis.lean}, \texttt{LeanTest1.lean}) enthalten keine
Mathlib-abhängigen Beweise; sie sind Platzhalter/Experimente.
\end{tcolorbox}

% ====================================================
\section{Bereits formalisierte Bausteine}
\label{sec:vorhanden}
% ====================================================

\subsection{2-adische Valuation (\texttt{padicValNat})}

\begin{tcolorbox}[vorhanden,
  title={Formalisiert: 2-adische Valuation in Mathlib}]
\textbf{Paket:} \Mathlib{NumberTheory.Padics.PadicVal.Defs} und
\Mathlib{NumberTheory.Padics.PadicVal.Basic}

\medskip
\textbf{Kerntheoreme:}
\begin{itemize}[nosep]
  \item \texttt{padicValNat p n} --- die $p$-adische Valuation $\vii(n)$ für $n:\N$
  \item \texttt{padicValNat.self} : \texttt{padicValNat p p = 1}
  \item \texttt{padicValNat\_pow} : \texttt{padicValNat p (p\^{}n) = n}
  \item \texttt{padicValNat.mul} : Multiplikativität $\vii(ab) = \vii(a)+\vii(b)$ (für $a,b\neq 0$)
  \item \texttt{padicValNat\_le\_nat\_log} : $\vii(n) \leq \lfloor\log_p n\rfloor$
  \item \texttt{padicValNat\_factorial} : Legendre-Satz $\vii(n!) = \sum_{i\geq 1}\lfloor n/p^i\rfloor$
  \item \texttt{padicValNat\_choose} : Kummer-Satz (Träger von $\binom{n}{k}$)
  \item \texttt{one\_le\_padicValNat\_of\_dvd} : $p\mid n \Rightarrow \vii(n)\geq 1$
\end{itemize}

\medskip
\textbf{Verbindung zum Dichte-Argument:} Das Residuenklassen-Argument
(C-Ketten-Startpunkte $\equiv 2^{k+1}-1 \pmod{2^{k+2}}$) ist über
\texttt{ZMod} und \texttt{Nat.dvd\_iff} zugänglich.
Die Zählung in arithmetischen Progressionen ist über
\texttt{Nat.card\_Icc\_mod} formalisierbar.
\end{tcolorbox}

\begin{tcolorbox}[lücke, title={Fehlt: Collatz-spezifisches Dichte-Lemma}]
Das Lemma
\[
  \mathrm{density}\bigl(\{n \in 2\N+1 \mid \vii(n+1) \geq k\}\bigr) = 2^{-k}
\]
als \emph{formalisiertes} Theorem über \texttt{Nat.density} existiert nicht in Mathlib.
Die Infrastruktur (Residuenklassen, padicValNat) ist vorhanden; das spezifische
Dichte-Lemma für C-Ketten muss als eigenes Lean-Lemma geführt werden.

\textbf{Verfügbare Brücke:}
\begin{lstlisting}[language=Lean4]
-- Residuenklassen-Charakterisierung (direkt beweisbar):
-- n + 1 ≡ 0 (mod 2^(k+1)) ↔ padicValNat 2 (n+1) ≥ k+1
-- Über: padicValNat_dvd_iff_of_ne_one
\end{lstlisting}
\end{tcolorbox}

\subsection{LTE-Valuation (Lifting the Exponent)}

\begin{tcolorbox}[teilweise, title={Teilweise vorhanden: Multiplikativität und Potenzen}]
Das Lifting-the-Exponent-Lemma für $p=2$,
\[
  \vii(3^j - 1) = \begin{cases} 1 & j \text{ ungerade,}\\ \vii(j)+2 & j \text{ gerade,}\end{cases}
\]
ist \emph{nicht} als fertiges Theorem in Mathlib. Jedoch sind vorhanden:
\begin{itemize}[nosep]
  \item \texttt{multiplicity.Nat.prime\_pow\_eq\_one} (Mathlib3-Äquivalent; in Mathlib4
        über \texttt{padicValNat\_pow})
  \item \texttt{padicValNat.mul} und Ordnungs-/Multiplizitäts-Lemmata für das
        Gerüst des LTE-Beweises
\end{itemize}

Der LTE-Beweis für $p=2$, $n=3^j-1$ muss selbst geführt werden.
\end{tcolorbox}

% ====================================================
\subsection{Geometrische Reihen und Potenzreihen}
% ====================================================

\begin{tcolorbox}[vorhanden,
  title={Formalisiert: Geometrische Reihen (inkl. $x=1/2$-Spezialfall)}]
\textbf{Paket:} \Mathlib{Analysis.SpecificLimits.Basic}

\medskip
\textbf{Kerntheoreme:}
\begin{itemize}[nosep]
  \item \texttt{tsum\_geometric\_of\_lt\_one} : $\sum_{n=0}^\infty r^n = (1-r)^{-1}$ für $0\leq r < 1$
  \item \texttt{hasSum\_geometric\_two'(a)} :
    $\sum_{n=0}^\infty a\cdot (1/2)^{n+1} = a$ \quad \textit{(direkt!)}
  \item \texttt{tsum\_geometric\_two'(a)} : entsprechender tsum-Wert
  \item \texttt{NNReal.hasSum\_geometric} : für $r:\texttt{NNReal}$ mit $r < 1$
  \item \texttt{NNReal.tsum\_geometric} : $\sum r^n = (1-r)^{-1}$ in $\texttt{NNReal}$
\end{itemize}

\textbf{Für das Dichte-Argument direkt verfügbar:}
\begin{lstlisting}[language=Lean4]
-- Konvergenz der gewichteten Summe: ∑ (k+1) * (1/2)^(k+1) = 2
-- Über: hasSum_geometric_two' mit (k+1)-Gewichtung
-- oder via tsum_geometric_of_lt_one + Ableitung der erzeugenden Funktion
-- (d/dx) [x/(1-x)] |_{x=1/2} = (1-x)^{-2} |_{x=1/2} = 4
-- Daher: ∑ (k+1) * (1/2)^(k+1) = 2 ✓
\end{lstlisting}

Die Schlüsselidentität $\sum_{k=0}^\infty (k+1)\cdot (1/2)^{k+1} = 2$
ist über \texttt{hasSum\_pow\_mul\_geometric\_of\_abs\_lt\_one} beweisbar.
\end{tcolorbox}

% ====================================================
\section{Stirling + Gamma in Mathlib}
\label{sec:stirling}
% ====================================================

\begin{tcolorbox}[vorhanden,
  title={Formalisiert: Stirling-Formel (vollständig)}]
\textbf{Paket:} \Mathlib{Analysis.SpecialFunctions.Stirling}

\medskip
Die Stirling-Formel $n! \sim \sqrt{2\pi n}\,(n/e)^n$ ist in Mathlib
\textbf{vollständig formalisiert}, einschließlich:

\begin{itemize}[nosep]
  \item \texttt{Stirling.stirlingSeq n} $:= \dfrac{n!\cdot 2^n}{(n/e)^n}$
    --- die normierte Stirling-Folge
  \item \texttt{Stirling.tendsto\_stirlingSeq\_sqrt\_pi} :
    $\texttt{stirlingSeq} \xrightarrow{n\to\infty} \sqrt{\pi}$
  \item \texttt{Stirling.factorial\_isEquivalent\_stirling} :
    asymptotische Äquivalenz $n! \sim \sqrt{2\pi n}(n/e)^n$
  \item \texttt{Stirling.le\_factorial\_stirling} :
    untere Schranke $\sqrt{2\pi n}(n/e)^n \leq n!$ für \emph{alle} $n$
  \item \texttt{Stirling.sqrt\_pi\_le\_stirlingSeq} :
    $\sqrt{\pi} \leq \texttt{stirlingSeq}\,n$ für alle $n > 0$ (effektiv!)
  \item \textbf{Neu (2026):} \texttt{Stirling.log\_stirlingSeq\_sdiff\_le} :
    Robbins-Schranke $\log(\texttt{stirlingSeq}\,n) - \log(\texttt{stirlingSeq}(n+1))
    \leq \tfrac{1}{12n(n+1)}$
\end{itemize}

\textbf{Wichtig für den Collatz-Kontext:}
Die Stirling-Formel erlaubt asymptotische Abschätzungen für $|B_{2n}|$
via $\frac{2(2n)!}{(2\pi)^{2n}} \sim \sqrt{\pi n}(n/(\pi e))^{2n}$.
\end{tcolorbox}

\begin{tcolorbox}[vorhanden, title={Formalisiert: Gamma-Funktion}]
\textbf{Paket:} \Mathlib{Analysis.SpecialFunctions.Gamma.Basic}

\begin{itemize}[nosep]
  \item \texttt{Real.Gamma} --- vollständige Definition
  \item \texttt{Real.Gamma\_nat\_eq\_factorial} :
    $\Gamma(n+1) = n!$ für $n:\N$
  \item \texttt{Real.Gamma\_pos\_of\_pos} : $\Gamma(x) > 0$ für $x > 0$
  \item \texttt{Complex.Gamma\_nat\_eq\_factorial} : komplexe Version
\end{itemize}
\end{tcolorbox}

\begin{tcolorbox}[lücke,
  title={Fehlt: Asymptotik der Bernoulli-Zahlen via Stirling}]
Die asymptotische Formel
\[
  |B_{2n}| \;\sim\; \frac{2\,(2n)!}{(2\pi)^{2n}}
\]
ist \emph{nicht} als fertiges Theorem in Mathlib. Die Bausteine
($\Gamma$, Stirling, $\zeta(2n)$-Formel) sind vorhanden; die Kombination
muss selbst geführt werden.
\end{tcolorbox}

% ====================================================
\section{Bernoulli-Zahlen und von-Staudt-Clausen}
\label{sec:bernoulli}
% ====================================================

\begin{tcolorbox}[vorhanden,
  title={Formalisiert: Bernoulli-Zahlen und von-Staudt-Clausen}]
\textbf{Paket:} \Mathlib{NumberTheory.Bernoulli}

\medskip
\textbf{Vollständig formalisierte Theoreme:}

\begin{itemize}[nosep]
  \item \texttt{bernoulli' n} --- Definition via Wohlordnungsinduktion:
    $B_n = 1 - \sum_{k<n}\binom{n}{k}\frac{B_k}{n+1-k}$
  \item \texttt{bernoulli'\_eq\_zero\_of\_odd} : $B_{2j+1} = 0$ für $j \geq 1$
  \item \texttt{Bernoulli.vonStaudt\_clausen (k)} :
    \[
      B_{2k} + \sum_{\substack{p\text{ prim}\\(p-1)\mid 2k}} \frac{1}{p}
      \;\in\; \Z
    \]
    \textbf{(von-Staudt-Clausen vollständig bewiesen!)}
  \item \texttt{sum\_range\_pow} --- Faulhaber-Formel:
    $\sum_{k=0}^{n-1} k^p = \sum_{i=0}^p B_i \binom{p+1}{i}\frac{n^{p+1-i}}{p+1}$
  \item \texttt{sum\_Ico\_pow} --- alternative Faulhaber-Form
  \item \texttt{bernoulliPowerSeries} --- erzeugende Funktion $\sum B_n t^n/n!$
\end{itemize}

\medskip
\textbf{Verbindung zu $\zeta(2n)$:}
Die Formel
$\zeta(2n) = (-1)^{n+1}\frac{B_{2n}(2\pi)^{2n}}{2(2n)!}$
ist als \texttt{riemannZeta\_two\_mul\_nat} in Mathlib formalisiert
(\Mathlib{NumberTheory.ZetaFunction}).

\textbf{Relevanz für Bernoulli-Normschalen $\mathcal{I}(n)$:}
Der von-Staudt-Clausen-Satz sichert die Ganzzahligkeit des Nenners von $B_{2n}$.
Die Normschalen-Formel
$\mathcal{I}(n) = \frac{2(2n)!\,\zeta(2n)}{(2\pi)^{2n}}\cdot D(n) \in \Z$
ist damit auf der Bernoulli-Seite durch Mathlib abgedeckt.
\end{tcolorbox}

\begin{tcolorbox}[teilweise, title={Teilweise vorhanden: Bernoulli-Polynome}]
\textbf{Paket:} \Mathlib{NumberTheory.BernoulliPolynomials}

Die Bernoulli-Polynome $B_n(x)$ sind definiert und grundlegende Eigenschaften
sind bewiesen. Die Verbindung zur Zeta-Funktion über die Hurwitz-Zeta-Funktion
ist teilweise vorhanden; volle asymptotische Entwicklungen fehlen.
\end{tcolorbox}

% ====================================================
\section{Vier-Quadrate-Satz und Quaternionen}
\label{sec:vierquadrate}
% ====================================================

\begin{tcolorbox}[vorhanden, title={Formalisiert: Lagrange-Vierquadratesatz}]
\textbf{Paket:} \Mathlib{NumberTheory.SumFourSquares}

\begin{itemize}[nosep]
  \item \texttt{Nat.sum\_four\_squares (n)} :
    $\forall n \in \N,\; \exists a\,b\,c\,d \in \Z:\; n = a^2+b^2+c^2+d^2$
    \quad \textbf{(vollständig!)}
  \item \texttt{euler\_four\_squares} : Euler-Vier-Quadrate-Identität
    $(a^2+b^2+c^2+d^2)(x^2+y^2+z^2+w^2) = \ldots$
  \item \texttt{Nat.euler\_four\_squares} : Version für $\N$
  \item \texttt{Nat.Prime.sum\_four\_squares} : Primzahlversion
\end{itemize}

\textbf{Relevanz für \texttt{collatz\_vier\_teilnormen.tex}:}
Theorem~\ref{sec:vierquadrate} (Lagrange) ist direkt verfügbar.
Die Formel $N(3q+1) = 9N(q)+6\,\mathrm{Re}(q)+1$ ist eine elementare
Rechnung; die Quaternionen-Algebra ist über \texttt{Mathlib.Algebra.Quaternion}
zugänglich.
\end{tcolorbox}

\begin{tcolorbox}[vorhanden, title={Formalisiert: Quaternionen-Algebra}]
\textbf{Paket:} \Mathlib{Algebra.Quaternion}

\begin{itemize}[nosep]
  \item \texttt{Quaternion} --- vollständige Typ-Definition
  \item \texttt{Quaternion.normSq} --- Norm $N(q) = a^2+b^2+c^2+d^2$
  \item \texttt{Quaternion.normSq\_mul} : $N(qr) = N(q)N(r)$
  \item \texttt{Quaternion.re} --- Realteilprojektion $\mathrm{Re}(q)$
  \item Hurwitz-Quaternionen: teilweise über \texttt{EvenOdd}-Klassifikation
\end{itemize}
\end{tcolorbox}

\begin{tcolorbox}[lücke, title={Fehlt: Jacobi-Vierquadratensatz (Zählformel)}]
Die Jacobi-Formel
\[
  r_4(n) = 8\sum_{\substack{d\mid n \\ 4\nmid d}} d
\]
als formalisiertes Theorem ist \emph{nicht} in Mathlib vorhanden.
(Nur die Existenz der Zerlegung ist bewiesen, nicht die Anzahl.)
\end{tcolorbox}

% ====================================================
\section{Ergodizität und Birkhoff-Mittelwertsatz}
\label{sec:ergodic}
% ====================================================

\begin{tcolorbox}[vorhanden, title={Formalisiert: Ergodizitäts-Infrastruktur}]
\textbf{Pakete:} \Mathlib{Dynamics.Ergodic.Ergodic},
\Mathlib{Dynamics.BirkhoffSum.Average},
\Mathlib{Analysis.InnerProductSpace.MeanErgodic}

\begin{itemize}[nosep]
  \item \texttt{Ergodic f μ} --- vollständige Definition: maßerhaltend + pr\"a-ergodisch
  \item \texttt{PreErgodic} : messbar-invariante Mengen sind fast leer oder voll
  \item \texttt{QuasiErgodic} : schwächere Version (quasi-maßerhaltend)
  \item \texttt{birkhoffAverage R f g n x} :
    $\frac{1}{n}\sum_{k=0}^{n-1} g(f^{[k]}(x))$ --- Birkhoff-Mittel
  \item \textbf{von-Neumann-Mittelwertsatz (vollständig):}
    \texttt{ContinuousLinearMap.tendsto\_birkhoffAverage\_orthogonalProjection}:
    Birkhoff-Mittel $\to$ Orthogonalprojektion auf Fixpunktunterraum
  \item \textbf{von-Neumann für Banachräume:}
    \texttt{tendsto\_birkhoffAverage} (allgemeine Version)
  \item \textbf{Gallagher's Ergodensatz:} formalisiert für den Einheitskreis
    (\Mathlib{NumberTheory.WellApproximable})
\end{itemize}
\end{tcolorbox}

\begin{tcolorbox}[lücke, title={Fehlt: Birkhoff-Punktgrenzwertsatz (Punktweise Version)}]
Der \emph{punktweise} Birkhoff-Ergodensatz
($\frac{1}{n}\sum_{k<n} g(f^k(x)) \to \int g\,d\mu$ f.\,ü.)
ist \textbf{noch nicht} in Mathlib formalisiert (Stand: Juni 2026).

Der von-Neumann-Mittelwertsatz ($L^2$-Konvergenz) ist vorhanden;
die punktweise Aussage (f.ü.-Konvergenz) fehlt.

\textbf{Für Collatz:} Ein Ergodizitätsbeweis für die Collatz-Abbildung
$\mathcal{U}$ auf ungeraden Zahlen würde den punktweisen Birkhoff-Satz benötigen.
\end{tcolorbox}

\begin{tcolorbox}[lücke, title={Fehlt: Kingman-Subadditivitätssatz}]
Der \textbf{Kingman-Unteradditivitätssatz} (subadditive ergodic theorem):
\[
  \frac{1}{n}\,g_n(x) \;\xrightarrow{n\to\infty}\; \inf_n \frac{1}{n}\int g_n\,d\mu
  \quad\text{f.\,ü.}
\]
für subadditive Folgen $(g_n)$ --- ist \textbf{nicht} in Mathlib.

\textbf{Relevanz:} Der Lyapunov-Exponent der Collatz-Abbildung als
$\Lambda = \lim_{n\to\infty}\frac{1}{n}\log|D\mathcal{U}^n|$ erfordert Kingman für
eine rigorose Existenzaussage.
\end{tcolorbox}

% ====================================================
\section{Direkt anwendbare Theoreme für das Dichte-Argument}
\label{sec:direkt}
% ====================================================

In \texttt{collatz\_lte\_dichte.tex} sind folgende Schlüsselberechnungen durchgeführt.
Hier zeigen wir, welche Lean-Beweise \emph{direkt} mit Mathlib machbar sind.

\begin{tcolorbox}[vorhanden, title={Beweisbar: Konvergenz $\sum(k+1)\cdot 2^{-(k+1)} = 2$}]
\begin{lstlisting}[language=Lean4]
-- Import: Mathlib.Analysis.SpecificLimits.Basic
-- Direkt via tsum der gewichteten geometrischen Reihe:
theorem sum_kp1_half_kp1 :
    ∑' k : ℕ, (k + 1 : ℝ) * (1/2)^(k+1) = 2 := by
  -- Via Ableitung: d/dx [x/(1-x)] = (1-x)^{-2}, bei x=1/2 Wert 4
  -- ∑ (k+1) x^(k+1) = x/(1-x)^2, x=1/2: 2
  -- Oder: hasSum_pow_mul_geometric_of_abs_lt_one mit n=1
  sorry -- Vollständig beweisbar mit Mathlib-Mitteln
\end{lstlisting}

\textbf{Mathlib-Theoreme dafür:}
\begin{itemize}[nosep]
  \item \texttt{hasSum\_pow\_mul\_geometric\_of\_abs\_lt\_one}
  \item \texttt{tsum\_geometric\_of\_lt\_one}
  \item \texttt{hasSum\_geometric\_two'} (für den $(1/2)$-Spezialfall)
\end{itemize}
\end{tcolorbox}

\begin{tcolorbox}[vorhanden, title={Beweisbar: Tail-Formel $\sum_{k=K}^\infty (k+1)\cdot 2^{-(k+1)} = (K+2)\cdot 2^{-K}$}]
\begin{lstlisting}[language=Lean4]
-- Via Differenz der Gesamtreihe minus endlicher Partialsumme:
-- Mathlib: tsum_eq_zero_add, sum_geometric_two
-- + Linearität von tsum
theorem extremal_tail (K : ℕ) :
    ∑' k : ℕ, (k + K + 1 : ℝ) * (1/2)^(k+K+1) = (K+2) * 2^(-K : ℤ) := by
  sorry -- Beweisbar mit tsum-Lemmas + Indexverschiebung
\end{lstlisting}
\end{tcolorbox}

\begin{tcolorbox}[vorhanden, title={Beweisbar: $\Lambda_A = 2\log_2 3 - 5 < 0$}]
Der Wert $\Lambda_A = 2(\log_2 3 - 1) - 3 \cdot 1 = 2\log_2 3 - 5$
ist eine geschlossene reelle Zahl. Die Negativität $2\log_2 3 < 5$
ist über \texttt{Real.log\_lt'} und numerische Schranken direkt beweisbar:
\begin{lstlisting}[language=Lean4]
-- In Mathlib: Real.log_two_pos, Nat.log_lt_of_lt_pow
-- Numerisch: log₂ 3 < 1.58497 → 2 * 1.58497 = 3.16994 < 5 ✓
theorem lambda_A_neg : 2 * Real.log 3 / Real.log 2 - 5 < 0 := by
  norm_num [Real.log_lt_log_iff]  -- mit norm_num-Erweiterung
\end{lstlisting}
\end{tcolorbox}

\begin{tcolorbox}[teilweise, title={Teilweise beweisbar: LTE-Valuation $\vii(3^j-1)$}]
\begin{lstlisting}[language=Lean4]
-- Ziel:
lemma lte_padic_3pow_sub_one (j : ℕ) :
    padicValNat 2 (3^j - 1) =
    if j % 2 = 1 then 1
    else padicValNat 2 j + 2 := by
  -- Beweisbar via:
  -- padicValNat.mul, Ord_2(3-1)=1, Ord_2(3+1)=2
  -- + Ordnungstheorie für 3 in (ℤ/2^k)^×
  -- KEIN fertiges Mathlib-Theorem; erfordert eigene Hilfs-Lemmata
  sorry
\end{lstlisting}
Das Gerüst ist via \texttt{padicValNat}, \texttt{ZMod.orderOf} und
\texttt{Nat.ord\_compl\_dvd} aufbaubar, aber nicht fertig formalisiert.
\end{tcolorbox}

% ====================================================
\section{Was noch fehlt}
\label{sec:luecken}
% ====================================================

\begin{tcolorbox}[lücke, title={Fehlt: Kingman-Subadditivitätssatz}]
\textbf{Benötigt für:} Lyapunov-Exponent-Existenz in \texttt{collatz\_lyapunov.tex}

\textbf{Status:} Nicht in Mathlib (Juni 2026). Der Satz erfordert
punktweise fast-überall-Konvergenz subadditiver Folgen --- eng verwandt
mit dem fehlenden Birkhoff-Punktgrenzwertsatz.

\textbf{Workaround:} Für spezifische (nicht-ergodische) Abschätzungen
kann man auf \texttt{liminf}-Lemmas und Feketes Lemma zurückgreifen.
\texttt{Real.tendsto\_of\_bddBelow\_of\_antitone} ist verfügbar.
\end{tcolorbox}

\begin{tcolorbox}[lücke, title={Fehlt: Birkhoff-Punktgrenzwertsatz}]
\textbf{Benötigt für:} $\Lambda = \lim_{N\to\infty}\frac{1}{N}\sum_{k<N}\log_2 F_k$
als f.\,ü.-Konvergenz

\textbf{Status:} Nicht in Mathlib. Der von-Neumann-Satz ($L^2$) ist da;
die $L^1$-pointwise Version fehlt. Dies ist ein bekanntes Mathlib-Desiderat.
\end{tcolorbox}

\begin{tcolorbox}[lücke, title={Fehlt: Collatz-spezifische Strukturlemmata}]
Die folgenden Lemmata sind vollständig neu (kein Mathlib-Äquivalent):

\begin{enumerate}[nosep, label=(\arabic*)]
  \item \texttt{cChainLength} --- Länge der C-Kette: muss definiert werden
  \item \texttt{density\_C\_chains} --- Dichte-Lemma (Abschnitt~\ref{sec:direkt})
  \item \texttt{lte\_worst\_forces\_E\_type} --- strukturelle Kompensation
  \item \texttt{block\_drift\_neg} --- negativer Block-Drift
  \item \texttt{collatz\_ergodic} --- Ergodizität der Collatz-Abbildung
\end{enumerate}
\end{tcolorbox}

\begin{tcolorbox}[lücke, title={Fehlt: Asymptotik der Bernoulli-Zahlen}]
\[
  |B_{2n}| \;\sim\; \frac{2\,(2n)!}{(2\pi)^{2n}}
  \quad (n\to\infty)
\]
Nicht in Mathlib. Bausteine ($\Gamma$, Stirling, $\zeta$) sind vorhanden;
die Zusammensetzung muss geführt werden.
\end{tcolorbox}

\begin{tcolorbox}[lücke, title={Fehlt: Jacobi-Vierquadratsatz (Zählformel $r_4$)}]
\[
  r_4(n) = 8\sum_{\substack{d\mid n\\ 4\nmid d}} d
\]
Die \emph{Existenz} einer Zerlegung ist formalisiert (\texttt{Nat.sum\_four\_squares});
die Zählung der Zerlegungsdarstellungen fehlt.
\end{tcolorbox}

% ====================================================
\section{Roadmap: 5--10 Lean-Beweise mit größtem Impact}
\label{sec:roadmap}
% ====================================================

Die folgende Priorisierung richtet sich nach dem Impact für den Collatz-Beweis
in den Dokumenten \texttt{collatz\_lte\_dichte.tex},
\texttt{collatz\_vier\_teilnormen.tex} und \texttt{collatz\_bernoulli\_schalen.tex}.

\begin{tcolorbox}[roadmap, title={Priorität 1: Dichte-Lemma für C-Ketten-Startpunkte}]
\textbf{Lemma (Impact: hoch --- Fundament des Dichte-Arguments)}
\begin{lstlisting}[language=Lean4]
-- Datei: Collatz/Density.lean
-- Mathlib-Importe: Mathlib.NumberTheory.Padics.PadicVal.Basic
--                  Mathlib.Data.Nat.Defs

def cChainStart (k : ℕ) : Set ℕ :=
  {n : ℕ | Odd n ∧ padicValNat 2 (n + 1) ≥ k + 1}

-- Ziel: Dichte = 2^(-(k+1)) in den ungeraden Zahlen
theorem density_C_chains (k : ℕ) (N : ℕ) (hN : N > 0) :
    #{n ∈ Finset.range N | n ∈ cChainStart k} =
    N / 2^(k+1) - N / 2^(k+2) := by
  -- Über: Residuenklasse n+1 ≡ 2^(k+1) (mod 2^(k+2))
  -- Schritt: padicValNat 2 (n+1) ≥ k+1 ↔ 2^(k+1) ∣ (n+1)
  -- + Nat.card_filter_range_mod
  sorry
\end{lstlisting}
\textbf{Mathlib-Basis:} \texttt{padicValNat\_dvd\_iff\_of\_ne\_one},
\texttt{Nat.card\_Ioc\_filter\_dvd}, arithmetische Progressionen.
\textbf{Schwierigkeit:} mittel (2--3 Wochen).
\end{tcolorbox}

\begin{tcolorbox}[roadmap, title={Priorität 2: Strukturelle Kompensation (LTE-Worst-Case)}]
\textbf{Lemma (Impact: hoch --- Zwei-Block-Mechanismus)}
\begin{lstlisting}[language=Lean4]
-- Datei: Collatz/StructuralComp.lean

-- Definition der EABC-Klassen via (n % 12):
def eabcClass (n : ℕ) : Fin 4 :=
  if n % 12 = 1 then 0  -- E
  else if n % 12 = 5 then 1  -- A
  else if n % 12 = 7 then 2  -- B
  else 3  -- C

-- Hauptlemma: Extremale Familie erzwingt E-Typ-Nachfolger
theorem lte_worst_forces_E_type (k r : ℕ) (hk : k ≥ 1) (hr : r ≥ 1) :
    let n0 := 2^(k+1) * 3^r - 1
    eabcClass (collatzU^[k+1] n0) = 0 ∨  -- E-Typ
    padicValNat 2 (3 * collatzU^[k] n0 + 1) ≥ 3 := by
  sorry
\end{lstlisting}
\textbf{Kernaussage:} Proposition~\ref{sec:lte} aus \texttt{collatz\_lte\_dichte.tex}
ist bewiesen; die Formalisierung erfordert arithmetische Modularität.
\textbf{Schwierigkeit:} mittel-hoch (4--6 Wochen, intensive Fallunterscheidungen).
\end{tcolorbox}

\begin{tcolorbox}[roadmap, title={Priorität 3: Block-Drift und negative Kontraktion}]
\textbf{Lemma (Impact: mittel --- zentrales Konvergenzkriterium)}
\begin{lstlisting}[language=Lean4]
-- Datei: Collatz/BlockDrift.lean

-- Lyapunov-Exponent im typischen Fall
theorem block_drift_neg :
    ∑' k : ℕ, (k + 1 : ℝ) * (1/2)^(k+1) *
    ((k+1) * Real.log 3 - 3) / Real.log 2
    < 0 := by
  -- Numerisch: Λ_A = 2 log₂ 3 - 5 ≈ -1.83 < 0
  -- Bausteine: tsum_geometric_two', Real.log_lt'
  sorry

-- Äquivalent:
theorem lambda_A_value :
    2 * Real.log 3 / Real.log 2 - 5 =
    ∑' k : ℕ, (k+1 : ℝ) * (1/2)^(k+1) *
    ((k+1) * (Real.log 3 / Real.log 2) - 3) := by
  sorry
\end{lstlisting}
\textbf{Mathlib-Basis:} \texttt{hasSum\_pow\_mul\_geometric\_of\_abs\_lt\_one},
\texttt{tsum\_mul\_left}, \texttt{Real.log\_lt\_log}.
\textbf{Schwierigkeit:} mittel (2--3 Wochen).
\end{tcolorbox}

\begin{tcolorbox}[roadmap, title={Priorität 4: von-Staudt-Clausen für $\mathcal{I}(n)$}]
\textbf{Lemma (Impact: mittel --- Ganzzahligkeit der Bernoulli-Normschalen)}
\begin{lstlisting}[language=Lean4]
-- Bereits in Mathlib! Nur Anwendung nötig:
-- Bernoulli.vonStaudt_clausen (k : ℕ) :
-- B_{2k} + ∑_{p prime, (p-1)∣2k} 1/p ∈ ℤ

-- Neue Aufgabe: Verbindung zu I(n)
-- I(n) = 2 * (2n)! * ζ(2n) / (2π)^{2n} * D(n)
-- Über riemannZeta_two_mul_nat + vonStaudt_clausen:
theorem I_n_is_integer (n : ℕ) (hn : n ≥ 1) :
    ∃ m : ℤ, (I_n n : ℝ) = m := by
  -- Kombination von:
  -- riemannZeta_two_mul_nat : ζ(2n) = ...B_{2n}...
  -- Bernoulli.vonStaudt_clausen
  sorry
\end{lstlisting}
\textbf{Mathlib-Basis:} \texttt{Bernoulli.vonStaudt\_clausen},
\texttt{riemannZeta\_two\_mul\_nat}. Weitgehend mechanisch.
\textbf{Schwierigkeit:} niedrig--mittel (1--2 Wochen nach Definitionen).
\end{tcolorbox}

\begin{tcolorbox}[roadmap, title={Priorität 5: Quaternionen-Norm-Identität}]
\textbf{Lemma (Impact: mittel --- Grundlage der Vier-Teilnormen-Analyse)}
\begin{lstlisting}[language=Lean4]
-- Datei: Collatz/QuaternionStep.lean
-- Import: Mathlib.Algebra.Quaternion

-- Direkte Berechnung (kein sorry nötig):
theorem quaternion_collatz_norm (a b c d : ℤ) :
    let q := (⟨a, b, c, d⟩ : ℍ[ℤ])
    Quaternion.normSq (3 * q + 1) =
    9 * Quaternion.normSq q + 6 * q.re + 1 := by
  simp [Quaternion.normSq, Quaternion.re]
  ring

-- Korollar: Asymmetrische Teilnorm-Transformation
theorem norm_collatz_step_re_only (a b c d : ℤ) :
    let q := (⟨a, b, c, d⟩ : ℍ[ℤ])
    Quaternion.normSq (3 * q + 1) = (3*a+1)^2 + 9*b^2 + 9*c^2 + 9*d^2 := by
  simp [Quaternion.normSq]
  ring
\end{lstlisting}
\textbf{Status:} Direkt beweisbar via \texttt{ring}; kein \texttt{sorry} nötig.
\textbf{Schwierigkeit:} niedrig (wenige Stunden).
\end{tcolorbox}

\begin{tcolorbox}[roadmap, title={Priorität 6: LTE-Valuation $\vii(3^j-1)$}]
\begin{lstlisting}[language=Lean4]
-- Datei: Collatz/LTE.lean
-- Import: Mathlib.NumberTheory.Padics.PadicVal.Basic

lemma padic_val_3pow_sub_one (j : ℕ) (hj : j ≥ 1) :
    padicValNat 2 (3^j - 1) =
    if j % 2 = 1 then 1 else padicValNat 2 j + 2 := by
  induction j with
  | zero => omega
  | succ n ih =>
    -- 3^(n+1) - 1 = (3-1)(3^n + ... + 1) = 2 * (...)
    -- Für n gerade: 3^n ≡ 1 mod 4, also 3^n + ... + 1 ≡ 2 mod 4
    -- Für n ungerade: ord₂(3^n - 1) = 1 (n ungerade)
    sorry
\end{lstlisting}
\textbf{Schwierigkeit:} mittel-hoch (Induktion + Ordnungstheorie).
\end{tcolorbox}

\begin{tcolorbox}[roadmap, title={Priorität 7: Ergodizitäts-Setup für Collatz}]
\begin{lstlisting}[language=Lean4]
-- Datei: Collatz/Ergodic.lean
-- Import: Mathlib.Dynamics.Ergodic.Ergodic
--         Mathlib.MeasureTheory.Measure.Probability

-- Definition der Collatz-Abbildung auf ungeraden ℕ:
noncomputable def collatzU : ℕ → ℕ :=
  fun n => (3 * n + 1) / 2^(padicValNat 2 (3*n+1))

-- Hypothetisches Ergodizitätstheorem (UNBEWIESEN):
-- theorem collatz_ergodic :
--     Ergodic collatzU collatzMeasure := by
--   sorry  -- Entspricht der Collatz-Vermutung!
-- ANMERKUNG: Ergodizität impliziert Collatz-Vermutung
\end{lstlisting}
\textbf{Schwierigkeit:} offen (entspricht im Wesentlichen der Collatz-Vermutung).
\end{tcolorbox}

\begin{tcolorbox}[roadmap, title={Priorität 8: Robbins-Schranke für Stirling-Anwendung}]
\begin{lstlisting}[language=Lean4]
-- Bereits in Mathlib (Priorität: sofort nutzbar!)
-- Stirling.log_stirlingSeq_sdiff_le (n : ℕ) :
-- log (stirlingSeq n) - log (stirlingSeq (n+1)) ≤ 1 / (12*n*(n+1))

-- Anwendung für Bernoulli-Normschalen:
-- |B_{2n}| kann über stirlingSeq nach oben abgeschätzt werden
-- ohne eigenen Stirling-Beweis!
theorem bernoulli_upper_via_stirling (n : ℕ) (hn : n ≥ 1) :
    |bernoulli (2*n)| ≤ 2 * Real.exp 1 * Stirling.stirlingSeq n^2 := by
  sorry -- Über: Bernoulli.vonStaudt_clausen + Stirling.le_factorial_stirling
\end{lstlisting}
\end{tcolorbox}

% ====================================================
\section{Zusammenfassung und Fazit}
\label{sec:fazit}
% ====================================================

\subsection{Status-Übersicht}

\begin{center}
\renewcommand{\arraystretch}{1.4}
\begin{tabular}{>{\raggedright}p{5.5cm}
                >{\centering}p{2.5cm}
                >{\raggedright\arraybackslash}p{5.0cm}}
\toprule
\textbf{Mathlib-Bereich} & \textbf{Status} & \textbf{Für Collatz nutzbar} \\
\midrule
\texttt{padicValNat}, Valuation-Infrastruktur
  & \textcolor{mathlibgrün}{\bfseries Vollständig}
  & Dichte-Argument-Fundament \\[2pt]
Geometrische Reihen ($x=1/2$)
  & \textcolor{mathlibgrün}{\bfseries Vollständig}
  & $\sum(k+1)\cdot 2^{-(k+1)} = 2$ direkt \\[2pt]
Stirling-Formel (inkl.\ Robbins)
  & \textcolor{mathlibgrün}{\bfseries Vollständig}
  & Bernoulli-Wachstumsschranken \\[2pt]
Bernoulli-Zahlen + von-Staudt-Clausen
  & \textcolor{mathlibgrün}{\bfseries Vollständig}
  & $\mathcal{I}(n) \in \Z$ direkt \\[2pt]
Lagrange-Vierquadratesatz
  & \textcolor{mathlibgrün}{\bfseries Vollständig}
  & Quaternionen-Norm-Analyse \\[2pt]
Quaternionen-Algebra (\texttt{normSq})
  & \textcolor{mathlibgrün}{\bfseries Vollständig}
  & Asymmetrie-Formel via \texttt{ring} \\[2pt]
Ergodizitäts-Definition + Birkhoff-Mittel
  & \textcolor{mathlibgrün}{\bfseries Vollständig}
  & von-Neumann-Mittelwertsatz \\[2pt]
$\zeta(2n)$-Formel via Bernoulli
  & \textcolor{mathlibgrün}{\bfseries Vollständig}
  & Normschalen-Ganzzahligkeit \\[2pt]
\midrule
Jacobi-Vierquadratsatz ($r_4$-Formel)
  & \textcolor{warnorange}{\bfseries Teilweise}
  & Existenz $\checkmark$; Zählung fehlt \\[2pt]
LTE-Valuation $\vii(3^j-1)$
  & \textcolor{warnorange}{\bfseries Teilweise}
  & Gerüst vorhanden; fertig fehlt \\[2pt]
Bernoulli-Asymptotik $|B_{2n}|$
  & \textcolor{warnorange}{\bfseries Teilweise}
  & Bausteine da; Komposition fehlt \\[2pt]
\midrule
Dichte-Lemma C-Ketten
  & \textcolor{lückenrot}{\bfseries Fehlt}
  & Neue Formalisierung nötig \\[2pt]
Birkhoff-Punktgrenzwertsatz (f.\,ü.)
  & \textcolor{lückenrot}{\bfseries Fehlt}
  & Mathlib-Desiderat \\[2pt]
Kingman-Subadditivitätssatz
  & \textcolor{lückenrot}{\bfseries Fehlt}
  & Mathlib-Desiderat \\[2pt]
Collatz-Ergodizität
  & \textcolor{lückenrot}{\bfseries Fehlt}
  & $\equiv$ Collatz-Vermutung \\[2pt]
Strukturelle Kompensation (Lean)
  & \textcolor{lückenrot}{\bfseries Fehlt}
  & Neue Formalisierung nötig \\[2pt]
\bottomrule
\end{tabular}
\end{center}

\subsection{Gesamtbewertung}

\begin{tcolorbox}[kernidee, title={Strategische Empfehlung}]
\textbf{Sofort nutzbar (keine neuen Lemmata):}
\begin{enumerate}[nosep, label=(\arabic*)]
  \item Quaternionen-Norm-Identität $N(3q+1) = 9n+6a+1$
    via \texttt{ring} in \Mathlib{Algebra.Quaternion}
  \item Von-Staudt-Clausen für $\mathcal{I}(n) \in \Z$
    via \texttt{Bernoulli.vonStaudt\_clausen}
  \item $\sum(k+1)\cdot 2^{-(k+1)} = 2$
    via \texttt{hasSum\_pow\_mul\_geometric\_of\_abs\_lt\_one}
  \item $\Lambda_A < 0$ numerisch
    via \texttt{norm\_num + Real.log\_lt'}
  \item Lagrange-Zerlegung $n = a^2+b^2+c^2+d^2$
    via \texttt{Nat.sum\_four\_squares}
\end{enumerate}

\medskip
\textbf{Nächster Schritt mit größtem Impact:}
Formalisierung des \textbf{Dichte-Lemmas} (Priorität~1), da es das
Fundament des gesamten Lyapunov-Dichte-Arguments bildet und direkt
mit vorhandener Mathlib-Infrastruktur (\texttt{padicValNat}) durchführbar ist.

\medskip
\textbf{Nicht kurzfristig formalisierbar:}
Kingman, punktweiser Birkhoff-Satz und Collatz-Ergodizität sind
mittelfristige Mathlib-Beiträge oder hängen an offenen mathematischen Fragen.
\end{tcolorbox}

\begin{tcolorbox}[warnung, title={Wichtige Einschränkung}]
Kein in Mathlib verfügbares Theorem beweist die Collatz-Vermutung.
Die hier gelisteten Formalisierungen stärken die Infrastruktur und
verifizieren Teilaussagen (Dichte, Strukturkompensation, Bernoulli-Ganzzahligkeit),
aber die DC-Lücke (uniforme Beschränktheit aller Orbits) bleibt offen.
\end{tcolorbox}

% ====================================================
\begin{thebibliography}{99}

\bibitem{mathlib4}
The Mathlib Community,
\textit{Mathlib4 -- Mathematics Library for Lean 4},
\url{https://leanprover-community.github.io/mathlib4_docs/},
Stand: Juni 2026.

\bibitem{stirling-mathlib}
\Mathlib{Analysis.SpecialFunctions.Stirling} ---
Stirling-Formel inkl.\ Robbins-Schranke.

\bibitem{bernoulli-mathlib}
\Mathlib{NumberTheory.Bernoulli} ---
Bernoulli-Zahlen, von-Staudt-Clausen, Faulhaber.

\bibitem{four-squares-mathlib}
\Mathlib{NumberTheory.SumFourSquares} ---
Lagrange-Vierquadratesatz.

\bibitem{padic-mathlib}
\Mathlib{NumberTheory.Padics.PadicVal.Basic} ---
$p$-adische Valuation, Legendre, Kummer.

\bibitem{ergodic-mathlib}
\Mathlib{Dynamics.Ergodic.Ergodic} ---
Ergodizitätsdefinition.

\bibitem{birkhoff-mathlib}
\Mathlib{Dynamics.BirkhoffSum.Average} ---
Birkhoff-Mittelwert.

\bibitem{mean-ergodic-mathlib}
\Mathlib{Analysis.InnerProductSpace.MeanErgodic} ---
von-Neumann-Mittelwertsatz.

\bibitem{hoffbauer-lte}
T.\ Hoffbauer,
\textit{2-adisches Dichte-Argument gegen die LTE-Barriere},
\texttt{collatz\_lte\_dichte.tex} (2026).

\bibitem{hoffbauer-vier}
T.\ Hoffbauer,
\textit{Vier-Teilnormen-Faktorisierung und Collatz-Dynamik},
\texttt{collatz\_vier\_teilnormen.tex} (2026).

\end{thebibliography}

\end{document}
